% !TEX root = ../main.tex
% !TEX program = lualatex
\documentclass[../main.tex]{subfiles}
\IfSubfilesClassLoaded{\externaldocument{\subfix{../build/main}}}{}
% =============================================================================
% CHAPTER 2: TECHNICAL AUTHORITY AND ENGINEERING AGENTS
% DESCRIPTION: Definitions, responsibilities, warranting, issue resolution, and EA roles.
% =============================================================================

\begin{document}
\IfSubfilesClassLoaded{\chapter{EDO Study Guide}}{}
%====================
% Contents here
%====================

% === TECHNICAL AUTHORITY AND ENGINEERING AGENTS ===
\section{Technical Authority and Engineering Agents}

\subsection{Summary}
\begin{description}
  \item[\textbf{Technical authority.}] \ac{ta} is an inherently governmental authority to establish,
  monitor, and approve technical standards, tools, and processes; for naval vessels, current statute also
  requires \ac{secnav} to designate a \ac{sta} for each naval vessel class~\autocite{SECNAV5400-15D,USC-10-8669b}.
  \item[\textbf{Programmatic authority.}] \acp{peo}, \acp{pae}, direct reporting program managers, and
  \acp{pm} own programmatic execution through cost, schedule, performance, contracting, and resourcing
  channels; technical authorities provide technically acceptable alternatives and risk judgments that
  cannot be treated as just another schedule preference~\autocite{SECNAV5400-15D,SECNAV5000-2G}.
  \item[\textbf{Waterfront issue resolution.}] Most waterfront problems start with the project team,
  \ac{lta}, or \ac{wfcheng}; unresolved or cross-domain issues move through the \ac{twh}, deputy warranting
  officer, NAVSEA Chief Engineer, and ultimately the \ac{syscom} commander as technical authority.
  \item[\textbf{CO underway caveat.}] A commanding officer retains command authority and responsibility for
  safe operation and mission accomplishment, but does not replace the warranted technical-authority chain.
\end{description}

\subsection{Definition of technical authority}
\begin{description}
  \item[\textbf{Definition.}] \ac{ta} is the authority, responsibility, and accountability to establish, monitor, and approve technical standards, tools, and processes in conformance with higher authority policy, requirements, architectures, and standards~\autocite{edo-2-1-5-technical-authority-2025}.
  \item[\textbf{Process focus.}] The exercise of \ac{ta} establishes and assures adherence to technical standards and policies, providing a range of technically acceptable alternatives with risk and value assessments~\autocite{edo-2-1-5-technical-authority-2025}.
\end{description}

\subsection{Ultimate technical authority}
\begin{description}
  \item[\textbf{Inherently governmental.}] Technical authority is an inherently governmental function assigned to \ac{syscom} commanders by \ac{secnav}~\autocite{SECNAV5400-15D,edo-2-1-5-technical-authority-2025}.
  \item[\textbf{Ships and most weapons systems.}] Commander, \ac{navsea} (COMNAVSEA) is the ultimate technical authority for ships and most weapons systems~\autocite{edo-2-1-5-technical-authority-2025}.
  \item[\textbf{C4ISR, cybersecurity, and IT.}] Commander, \ac{navwar} (COMNAVWAR) is the ultimate technical authority for \ac{c4isr}, cybersecurity, and \ac{it} systems~\autocite{edo-2-1-5-technical-authority-2025}.
  \item[\textbf{Policy basis.}] Title 10 and SECNAVINST 5400.15 series underpin the technical authority structure~\autocite{USC-10-8669b,SECNAV5400-15D,edo-2-1-5-technical-authority-2025}.
\end{description}

\subsection{Responsibilities of technical warrant holders}
Responsibilities of \acp{twh} include the following technical competencies~\autocite{edo-2-1-5-technical-authority-2025}.
\begin{description}
  \item[\textbf{Set technical standards.}] Create standards and tailor them to specific program applications while balancing requirements with intended missions.
  \item[\textbf{Technical area expertise.}] Provide expert testimony (risk assessments, technology readiness evaluations, critiques, mishap investigations, trouble reports, test plans, and test evaluation reviews).
  \item[\textbf{Safe and reliable operations.}] Ensure safety and reliability are addressed in technical documentation and act as certification authority when required.
  \item[\textbf{Effective systems engineering.}] Ensure engineering and technical products meet Navy needs, considering full life-cycle and system-of-systems relationships.
  \item[\textbf{Unbiased technical judgment.}] Provide leadership and accountability for engineering decisions; promote communication so technical requirements, issues, and risks are identified early.
  \item[\textbf{Stewardship.}] Forecast, shape, and develop the capability and capacity of engineering support networks.
  \item[\textbf{Accountability and integrity.}] Balance speed to decision, expertise of empowered organizations, and the need to elevate significant issues.
\end{description}

\subsection{Technical authority interfaces}
\begin{description}
  \item[\textbf{Risk discussion.}] Technical decisions are part of the cost, schedule, and performance risk discussion~\autocite{edo-2-1-5-technical-authority-2025}.
  \item[\textbf{Stakeholder balance.}] Stakeholders (\ac{pmo}, \ac{nsy}, Fleet) balance risk with technical requirements~\autocite{edo-2-1-5-technical-authority-2025}.
  \item[\textbf{Recommendations.}] Technical authorities assess risk and provide recommendations and acceptable alternatives to programmatic authorities when operational or programmatic challenges arise~\autocite{edo-2-1-5-technical-authority-2025}.
\end{description}

\subsection{Need for technical authority}
\subsubsection{USS Dolphin (AGSS-555) case}
\begin{description}
  \item[\textbf{Incident.}] On 21 May 2002, USS Dolphin was operating 100 miles off San Diego when the ship began flooding while surfaced to recharge the battery; fires broke out and most of the crew prepared to abandon ship, but two crew members prevented total loss~\autocite{edo-2-1-5-technical-authority-2025}.
  \item[\textbf{Technical risk.}] A black D-shaped sponge gasket was substituted for the specified red silicone without technical authority authorization; the Mk 54 shield side door opened in heavy seas, water entered the submarine, flooded spaces, and power was lost for de-watering~\autocite{edo-2-1-5-technical-authority-2025}.
\end{description}

\subsubsection{Why technical authority is required}
\begin{description}
  \item[\textbf{Disciplined processes.}] Develop and employ consistent, disciplined, collaborative engineering processes that provide safe, reliable, effective, integrated, timely, and affordable products for the Navy~\autocite{edo-2-1-5-technical-authority-2025}.
  \item[\textbf{Standards and tools.}] Establish, monitor, and approve standards, tools, and processes to provide risk-based solutions to the Fleet~\autocite{edo-2-1-5-technical-authority-2025}.
  \item[\textbf{Core process oversight.}] Oversee technology development, systems engineering, cost estimating, and test and evaluation~\autocite{edo-2-1-5-technical-authority-2025}.
  \item[\textbf{Technical infrastructure.}] Operate and sustain the most efficient technical infrastructure to support acquisition, fielding, and in-service support~\autocite{edo-2-1-5-technical-authority-2025}.
  \item[\textbf{Senior support.}] Support \ac{asnrdanda} and \ac{cno} analysis of mission areas, systems, and requirements~\autocite{edo-2-1-5-technical-authority-2025}.
\end{description}

\subsection{Technical authorities and warrant holders}
\subsubsection{Technical authorities}
\begin{description}
  \item[\textbf{Chief Systems Engineer (CSE).}] Facilitates, collaborates, and coordinates with Technical Domain Managers within their warranted area to ensure proper oversight of technical authority (for example, ballistic missile defense, \ac{c4isr}, space)~\autocite{edo-2-1-5-technical-authority-2025}.
  \item[\textbf{Technical Domain Manager.}] Provides deep technical expertise to programs and activities (marine engineering, structures, fluid mechanics)~\autocite{edo-2-1-5-technical-authority-2025}.
  \item[\textbf{Senior Technical Authority (STA).}] Certifies sufficient risk mitigation and maturity for critical systems to proceed past Milestone B. Title 10 USC \S 8669b requires \ac{secnav} to designate an \ac{sta} for each naval vessel class and, as amended for the current \ac{pae} construct, requires the \ac{sta} to report directly to the relevant \ac{pae}; the coursebook also notes that \acp{sta} are dual-hatted as \ac{cse} Deputy Warranting Officers~\autocite{USC-10-8669b,edo-2-1-5-technical-authority-2025}.
\end{description}

\subsubsection{Technical warrant holders}\label{sec:twhs}
\begin{description}
  \item[\textbf{Ship Design Manager (SDM).}] Manages systems engineering and design efforts for assigned platforms (in-service, new construction, future design)~\autocite{edo-2-1-5-technical-authority-2025}.
  \item[\textbf{Systems Integration Manager (SIM).}] Assists \acp{sdm} in systems engineering integration of complex warfare systems into platforms (for example, ASW or DDG 51/CG 47 combat systems)~\autocite{edo-2-1-5-technical-authority-2025}.
  \item[\textbf{Cost Engineering Manager (CEM).}] Ensures independent cost engineering and estimating for Navy programs (CVN, submarines, surface platforms)~\autocite{edo-2-1-5-technical-authority-2025}.
  \item[\textbf{Technical Area Expert (TAE).}] Navy expert in the assigned product technical area (propulsion, missiles, shock, combat systems)~\autocite{edo-2-1-5-technical-authority-2025}.
  \item[\textbf{Technical Process Owner (TPO).}] Provides definition and documentation for assigned technical processes~\autocite{edo-2-1-5-technical-authority-2025}.
  \item[\textbf{Waterfront Chief Engineer (WFCHENG).}] Leads and focuses technical efforts of \acp{syscom} from the waterfront (\ac{nsy}, regional maintenance centers, \acp{supship}, \ac{navwar})~\autocite{edo-2-1-5-technical-authority-2025}.
\end{description}

\subsection{NAVSEA warranting structure (overview)}
\begin{description}
  \item[\textbf{Leadership.}] COMNAVSEA; NAVSEA Chief Engineer (SEA 05); Deputy Warranting Officers~\autocite{edo-2-1-5-technical-authority-2025}.
  \item[\textbf{Deputy Warranting Officer domains.}] Cost Engineering and Industrial Analysis (SEA 05C); Explosive Ordnance Engineering (SEA 05E); Warfare Systems Engineering, Littoral and Mine Warfare (SEA 05M); Surface Ship Design and Systems Engineering (SEA 05D); Integrated Warfare Systems Engineering (SEA 05H); Undersea Warfare Systems Engineering (SEA 05N); Ship Integrity and Performance Engineering (SEA 05P); Readiness and Logistics (SEA 05R); Submarine and Submersible Design and Systems Engineering (SEA 05U); Aircraft Carrier Design and Systems Engineering (SEA 05V); Surface Warfare Systems Engineering (SEA 05W); Marine Engineering (SEA 05Z)~\autocite{edo-2-1-5-technical-authority-2025}.
  \item[\textbf{Warrant holders.}] See Section~\ref{sec:twhs}~\autocite{edo-2-1-5-technical-authority-2025}.
  \item[\textbf{Engineering agents.}] Engineering agents execute delegated technical tasks in support of warrant holders~\autocite{edo-2-1-5-technical-authority-2025}.
\end{description}

\subsection{Resolving waterfront technical issues}
\begin{description}
  \item[\textbf{TA Principles.}] Technical authority is exercised independently of programmatic cost, schedule, and performance constraints. It provides safe, reliable, and standardized solutions, using risk-informed engineering judgments to support the fleet.
  \item[\textbf{DFS.}] When a repair or modification departs from specifications, drawings, or standards, a \ac{dfs} must be submitted. The \ac{dfs} is routed through the \ac{lta} for technical review and disposition, and is tracked in the ship's \ac{csmp}~\autocite{edo-2-1-5-technical-authority-2025}.
  \item[\textbf{Escalation Chain for Repairs.}] If a waterfront technical issue or \ac{dfs} cannot be resolved locally, or exceeds local authority limits, it is elevated through the following chain:
    \begin{enumerate}
      \item \textbf{Waterfront Engineer / Project Team}: Identifies the deviation or issue.
      \item \textbf{Local Technical Authority (LTA) / Waterfront Chief Engineer (WFCHENG)}: Adjudicates within delegated local limits (e.g., at the Shipyard, RMC, or SUPSHIP).
      \item \textbf{Technical Warrant Holder (TWH)} (e.g., Ship Design Manager (SDM) or Technical Area Expert (TAE) at NAVSEA 05): Resolves issues crossing domain, design, or platform boundaries.
      \item \textbf{Deputy Warranting Officer (DWO)}: Reviews issues that cut across domains or have significant program impact.
      \item \textbf{NAVSEA Chief Engineer (CHENG, SEA 05)}: Resolves high-level disputes and sets overall engineering policy.
      \item \textbf{Commander, NAVSEA (COMNAVSEA)}: Serves as the ultimate technical authority for ships and weapons systems.
    \end{enumerate}
  \item[\textbf{Programmatic vs. Technical Disputes.}] Program Managers (\acp{pm}/PEOs) possess programmatic authority (cost, schedule, performance), while SYSCOMs/TWHs possess technical authority (safety, standards). A PM cannot override a technical authority's disapproval of a deviation. Disagreements must be elevated through the respective programmatic and technical chains, with \acp{pae}, \acp{peo}, the \ac{syscom} commander, and \ac{asnrdanda} involved as the issue level and decision authority require~\autocite{SECNAV5400-15D,SECNAV5000-2G}.
  \item[\textbf{Get to yes safely.}] Avoid saying no when possible, avoid analysis paralysis, and use engineering judgment when data is not available; clearly explain why an operational restriction is needed~\autocite{edo-2-1-5-technical-authority-2025}.
  \item[\textbf{Underway authority.}] The ship's \ac{co} retains command authority and responsibility for safe operation, mission accomplishment, and survivability while underway; emergent technical decisions should be coordinated with the appropriate technical authority as practicable, with deviations documented and adjudicated through the \ac{ta} chain~\autocite{SECNAV5400-15D,edo-2-1-5-technical-authority-2025}.
\end{description}

\subsection{Approving authorities}
\begin{description}
  \item[\textbf{Policy chain.}] U.S.C. Title 10, SECNAVINST 5430.7 series, and SECNAVINST 5400.15 series define assignment of responsibilities and authorities~\autocite{edo-2-1-5-technical-authority-2025}.
  \item[\textbf{SYSCOM role.}] \acp{syscom} provide in-service support to \acp{peo} and direct reporting program managers; they serve as technical authority and operational safety and certification authorities for assigned areas~\autocite{edo-2-1-5-technical-authority-2025}.
  \item[\textbf{Programmatic authority.}] \acp{peo} and direct reporting program managers exercise programmatic authority for assigned programs on behalf of the \ac{nae}~\autocite{edo-2-1-5-technical-authority-2025}.
  \item[\textbf{Sustainment oversight.}] \ac{asnrdanda} is responsible for overall supervision of sustainment; \ac{cno} is responsible for manning, training, and equipping forces, including sustainment resourcing~\autocite{edo-2-1-5-technical-authority-2025}.
\end{description}

\subsection{Engineering Agents (EAs)}
\begin{description}
  \item[\textbf{Definition.}] \acp{ea} are organizations (government or private contractor) with responsibilities delegated by assigning documents and tasking for functions within their assigned technical area, system, or mission~\autocite{edo-2-1-5-technical-authority-2025}.
  \item[\textbf{Lifecycle support.}] \acp{ea} provide technical services to \acp{twh}, program managers, and Fleet customers throughout a system's life-cycle~\autocite{edo-2-1-5-technical-authority-2025}.
\end{description}

\subsection{Types of Engineering Agents}
\begin{description}
  \item[\textbf{Five types.}] Technical Direction Agent (TDA), Design Agent (DA), System Integration Agent (SIA), Acquisition Engineering Agent (AEA), and \ac{isea}~\autocite{edo-2-1-5-technical-authority-2025}.
  \item[\textbf{Lifecycle coverage.}] Engineering agents exercise acquisition responsibilities across the acquisition life-cycle~\autocite{edo-2-1-5-technical-authority-2025}.
\end{description}

\subsection{Technical Direction Agent (TDA)}
\begin{description}
  \item[\textbf{Responsibilities.}] Develop and analyze engineering concepts; develop system specifications; evaluate design agent proposals; assist in developing the \ac{temp}~\autocite{edo-2-1-5-technical-authority-2025}.
  \item[\textbf{Purpose.}] Assist in establishing initial program concepts, performance specifications, \ac{rdte}, and alternative technical approaches~\autocite{edo-2-1-5-technical-authority-2025}.
\end{description}

\subsection{Design Agent (DA)}
\begin{description}
  \item[\textbf{Responsibilities.}] Develop initial and complete engineering design and prototypes; prepare technical data; prepare maintenance support documentation; review engineering change proposals and waivers or deviations for impact on design, performance, safety, producibility, maintainability, life-cycle cost, training, and interfaces~\autocite{edo-2-1-5-technical-authority-2025}.
  \item[\textbf{Purpose.}] Translate performance specifications into engineering specifications for the ship and its systems~\autocite{edo-2-1-5-technical-authority-2025}.
\end{description}

\subsection{System Integration Agent (SIA)}
\begin{description}
  \item[\textbf{Responsibilities.}] Ensure equipment and software compatibility; advise the system manager of interface problems; coordinate with other system managers; support \ac{tda} in \ac{te} as required; review engineering change proposals for impact on interfaces~\autocite{edo-2-1-5-technical-authority-2025}.
  \item[\textbf{Purpose.}] Ensure compatibility of all elements that make up a system and its interface requirements~\autocite{edo-2-1-5-technical-authority-2025}.
\end{description}

\subsection{Acquisition Engineering Agent (AEA)}
\begin{description}
  \item[\textbf{Responsibilities.}] Review technical manuals and \acp{mrc} for accuracy; analyze and report production costs and issues; assist in developing maintenance concepts and criteria for all levels of maintenance~\autocite{edo-2-1-5-technical-authority-2025}.
  \item[\textbf{Purpose.}] Provide logistics support prior to production or major modification improvement programs~\autocite{edo-2-1-5-technical-authority-2025}.
\end{description}

\subsection{In-Service Engineering Agent (ISEA)}
\begin{description}
  \item[\textbf{Responsibilities.}] Support and provide life-cycle maintenance; recommend corrections to deficiencies in all areas of life-cycle support; manage configuration of equipment; provide engineering assistance to the Fleet~\autocite{edo-2-1-5-technical-authority-2025}.
  \item[\textbf{Purpose.}] Delegated functions for overall engineering, test, maintenance, and logistics of in-service equipment~\autocite{edo-2-1-5-technical-authority-2025}.
\end{description}

%====================
% End of file
%====================
\IfSubfilesClassLoaded{
  \RenewDocumentCommand{\entryname}{}{\textbf{\color{Modern} Acronym}}
  \RenewDocumentCommand{\descriptionname}{}{\textbf{\color{Modern} Definition}}
  \printnoidxglossary[
    type=\acronymtype,
    title=Acronyms,
    style=long-booktabs
  ]}{}
\end{document}
